% ===== 原文 pp.12–14（PDF pp.13–15）=====
% [原文 p.12]
例如，对任意 Riemann 流形 $\cB$，靶为 $\cB\times S^1_R$ 与 $\cB\times S^1_{1/R}$ 的 sigma 模型
是等价的。更一般地，让我们考虑靶空间在拓扑上是空间 $\cB$ 与 $n$-环面 $T^n=(S^1)^n$ 的乘积
$\cB\times T^n$、但度规与 $B$-场具有混合分量的 sigma 模型。只要度规与 $B$-场在环面方向的平移下
不变，路径积分论证在此情形也适用。设 $x^i$ 是 $T^n$ 的坐标，周期为 $2\pi$；设 $x^a$ 是 $\cB$ 的
任意局部坐标。令 $E=g+2\pi B$，并设 $E^{ij}$ 是 $E_{ij}=g_{ij}+2\pi B_{ij}$ 的逆矩阵。设
$\tilde x_i$ 是对偶环面 $\tT^n$ 的坐标，周期为 $2\pi$。则 T-对偶理论是靶空间为 $\cB\times\tT^n$、
且 $\tE=\tg+2\pi\tB$ 由下式给出的 sigma 模型：
\begin{equation}
  \tE^{ij}=E^{ij},\qquad
  \tE^{i}_{b}=E^{ij}E_{jb},\qquad
  \tE_{a}^{j}=-E_{ai}E^{ij},\qquad
  \tE_{ab}=E_{ab}-E_{ai}E^{ij}E_{jb}. \tag{3.9}
\end{equation}

\subsection{T-对偶即镜像对称}

让我们考虑靶空间为圆柱 $\R\times S^1_R$ 的二维 $(2,2)$ 超对称非线性 sigma 模型；该圆柱是平坦
Kähler 流形，度规为 $\dd s^2=|\dd z|^2$，其中 $z$ 是周期为 $z\equiv z+2\pi R\ii$ 的复坐标。场变量是
映射 $\phi:\Sigma\to\C/2\pi R\ii\Z$ 与一个 Dirac 费米子 $\psi$，拉氏量为
\begin{equation}
  \cL=\frac{1}{2}|\partial_0\phi|^2-\frac{1}{2}|\partial_1\phi|^2
  +\frac{\ii}{2}\overline{\psi}_-(\partial_0+\partial_1)\psi_-
  +\frac{\ii}{2}\overline{\psi}_+(\partial_0-\partial_1)\psi_+. \tag{3.10}
\end{equation}
该系统是靶为 $S^1_R$ 的 sigma 模型与另一个由 $\phi$ 的实部（一个普通标量场）及 Dirac 费米子 $\psi$
组成的系统的乘积。对 $S^1_R$ 施加 T-对偶，得到靶为 $S^1_{1/R}$ 的 sigma 模型与同一个 $\Re(\phi)$
和 $\psi$ 的系统的乘积。换言之，即靶为 $\R\times S^1_{1/R}$ 的 sigma 模型与一个 Dirac 费米子的乘积。
对偶系统既与原系统等价，就也必须具有 $(2,2)$ 超对称。那么，场变量的超对称变换是什么呢？

我们回顾原系统中的超对称变换（见 (2.15) 与 (2.16)）：
\begin{equation}
  \delta\phi=\epsilon_+\psi_--\epsilon_-\psi_+, \tag{3.11}
\end{equation}
\begin{equation}
  \delta\psi_\pm=\pm\ii\overline{\epsilon}_\mp(\partial_0\pm\partial_1)\phi. \tag{3.12}
\end{equation}
注意 $\phi$ 是 B-手征的，
\begin{equation}
  \overline{Q}_+\phi=\overline{Q}_-\phi=0. \tag{3.13}
\end{equation}

% [原文 p.13]
让我们写 $\phi=\Re(\phi)+\ii X$，其中 $X$ 是到 $S^1_R=\R/2\pi R\Z$ 的映射。如前面那样把 $X$ 的
T-对偶变量记作 $\tX$，我们可以把对偶复变量写成 $\tphi=\Re(\phi)+\ii\tX$。回顾关系 (3.5)，我们看到
$\phi$ 与 $\tphi$ 由下式相联系：
\begin{equation}
  (\partial_0+\partial_1)\phi=(\partial_0+\partial_1)\tphi,\qquad
  (\partial_0-\partial_1)\overline{\phi}=(\partial_0-\partial_1)\tphi. \tag{3.14}
\end{equation}
应用 (3.11) 并利用运动方程 $(\partial_0\pm\partial_1)\psi_\mp=0$，我们得到
$\delta(\partial_0\pm\partial_1)\tphi=(\partial_0\pm\partial_1)(-\overline{\epsilon}_+\overline{\psi}_--\epsilon_-\psi_+)$。
积分之，得到对偶系统中的超对称变换：
\begin{equation}
  \delta\tphi=-\overline{\epsilon}_+\overline{\psi}_--\epsilon_-\psi_+, \tag{3.15}
\end{equation}
\begin{equation}
  \delta\psi_+=\ii\overline{\epsilon}_-(\partial_0+\partial_1)\tphi,\qquad
  \delta\overline{\psi}_-=\ii\epsilon_+(\partial_0-\partial_1)\tphi. \tag{3.16}
\end{equation}
我们看到，它并不取 (2.15) 与 (2.16)（或 (3.11) 与 (3.12)）那样的标准形式。特别地，$\tphi$ 不是
B-手征的，而是 A-手征的，
\begin{equation}
  \overline{Q}_+\tphi=Q_-\tphi=0. \tag{3.17}
\end{equation}
然而，若我们把变分参数改名为 $-\overline{\epsilon}_+\to\epsilon'_+$ 与
$\epsilon_-\to\epsilon'_-$（从而 $-\epsilon_+\to\overline{\epsilon}'_+$ 且
$\overline{\epsilon}_-\to\overline{\epsilon}'_-$），并把超荷改名为
\begin{equation}
  Q_+\to Q'_+,\quad Q_-\to\overline{Q}'_-,\quad
  \overline{Q}_+\to\overline{Q}'_+,\quad \overline{Q}_-\to Q'_-,
  \tag{3.18}
\end{equation}
则由 $\delta=\ii\epsilon'_+Q'_--\ii\epsilon'_-Q'_+-\ii\overline{\epsilon}'_+\overline{Q}'_-+\ii\overline{\epsilon}'_-\overline{Q}'_+$
给出的变换取标准形式。换言之，靶为 $\R\times S^1_R$ 与 $\R\times S^1_{1/R}$ 的 $(2,2)$ 超对称 sigma
模型之间存在一个等价，在该等价下超荷按 (3.18) 映射。也就是说，它们互为镜像。

让我们看看 D-膜在此 T-对偶下如何变换。首先，考虑缠绕在圆 $S^1_R$ 上、在 $\Re(\phi)$ 中处于固定
位置、并支撑一个平坦 $U(1)$ 规范势的 D1-膜。这是圆柱的拉格朗日子流形，因而是一个 A-膜。T-对偶把
它映为这样的 D-膜：它在 $\Re(\phi)$ 中处于相同位置，而在对偶圆 $S^1_{1/R}$ 中处于对应于该规范势
和乐的位置。也就是说，它是对偶圆柱某点处的 D0-膜。这是复子流形，因而是一个 B-膜。其次，考虑 D1-膜
（同样是 A-膜），它在圆 $S^1_R$ 中处于固定位置，但沿 $\Re(\phi)$ 方向延伸。T-对偶把它映为在
$S^1_{1/R}$ 与 $\Re(\phi)$ 两个方向上都延伸的 D2-膜，其中 $S^1_{1/R}$ 中的和乐对应于原 D1-膜在
$S^1_R$ 中的位置。这是一个 B-膜。我们看到，原模型中的 A-膜被映为对偶中的 B-膜。这确实是镜像对称
的一个性质。

当 $\Re(\phi)$ 也是周期变量时，同样的结论成立——靶为 $S^1_{R'}\times S^1_R$ 与
$S^1_{R'}\times S^1_{1/R}$ 的 $(2,2)$ 超对称 sigma 模型互为镜像。我们
% [原文 p.14]
也可以考虑具有一般（平坦）度规与（平坦）$B$-场的二环面 $S^1\times S^1$ 上的 sigma 模型。设 $x^1$ 与
$x^2$ 是第一个与第二个圆的坐标，周期均为 $2\pi$。一般的度规与 $B$-场可写成
\begin{equation}
  \dd s^2=R^2\left|\dd x^1+\tau\,\dd x^2\right|^2,\qquad
  B=B_{12}\,\dd x^1\wedge\dd x^2, \tag{3.19}
\end{equation}
其中 $\tau=\tau_1+\ii\tau_2$ 位于上半平面，$\tau_2>0$。此 $\tau$ 决定环面的复结构，而复化 Kähler
类可参数化为
\begin{equation}
  \omega-\ii B=\frac{1}{2\pi\ii}\rho\,\dd x^1\wedge\dd x^2
  \quad\text{其中}\quad\rho=2\pi B_{12}+\ii R^2\tau_2. \tag{3.20}
\end{equation}
现在让我们对 $x^1$ 圆施加 T-对偶，$x^2$ 保持不变。T-对偶理论的度规与 $B$-场可从 (3.9) 读出。把
T-对偶坐标记作 $\tx_1$（同样具有周期性 $2\pi$），它们为
\begin{equation}
  \dd\widetilde{s}^{\,2}=\frac{1}{R^2}\left|\dd\tx_1+\left(2\pi B_{12}+\ii R^2\tau_2\right)\dd x^2\right|^2,
  \qquad
  \tB=\frac{\tau_1}{2\pi}\,\dd\tx_1\wedge\dd x^2. \tag{3.21}
\end{equation}
特别地，对偶理论的 $\tau$ 与 $\rho$ 参数为
\begin{equation}
  \tilde\tau=2\pi B_{12}+\ii R^2\tau_2=\rho,\qquad
  \tilde\rho=2\pi\frac{\tau_1}{2\pi}+\ii\frac{1}{R^2}R^2\tau_2=\tau. \tag{3.22}
\end{equation}
复结构与复化 Kähler 参数在对偶模型之间确实被交换了。这是镜像对称的另一个性质。

这可以推广到更高维。让我们考虑靶空间为 $X=\cB\times T^n$、带有度规 $g$ 与 $B$-场 $B$（可以有
混合分量）的 sigma 模型。加上费米变量，我们可以把它变成超对称 sigma 模型；只要 $X$ 容许一个使
$g$ 为 Kähler 的复结构、且 $B$-场是平坦的（$\dd B=0$），它就具有 $(2,2)$ 超对称。该超对称模型
具有到 $\tX=\cB\times\tT^n$ 的 T-对偶，其中度规 $\tg$ 与 $B$-场 $\tB$ 由与玻色情形 (3.9) 相同的
公式给出。于是可以证明，在如下条件下该 T-对偶是镜像对称：$\cB$ 每点处的 $T^n$ 都是 $X$ 的拉格朗日
子流形，且 $B$-场在其上为零，即 $B|_{T^n}=0$。（在此情形，我们将说 $X\to\cB$ 是一个\emph{拉格朗日
环面纤维化}。）注意，这不过是：缠绕在 $\cB$ 某点处的 $T^n$ 上、支撑平坦 $U(1)$ 规范势的 Dn-膜是一个
A-膜。这与镜像对称交换 A-膜和 B-膜这一事实完全一致：正如我们所看到的，这样的 Dn-膜在 T-对偶下被映为
$\tX$ 某点处的 D0-膜。由于后者是 B-膜，若该 T-对偶是镜像对称，则原来的 Dn-膜必是 A-膜。还要注意，
对偶空间 $\tX$ 可以重构为原模型中缠绕在 $T^n$ 上的 Dn-膜的模空间。

% [Tutor Note: p.12 式 (3.9) 的指标位置是关键，不能凭文字层（它会抽成 $eEij$ 之类）抄写：环面方向
%  的指标记在上方（$\tE^{ij}=E^{ij}$）、底空间方向记在下方（$\tE_{ab}$）、混合分量两种指标各一
%  （$\tE^{i}_{b}$、$\tE_{a}^{j}$）。可用 $n=1$、$B=0$ 核对：此时 $E_{11}=g_{11}=R^2$，故
%  $E^{11}=1/R^2$，这正是单圆 T-对偶后度规 $\dd\tilde s^2=R^{-2}(\dd\tilde x_1)^2$ 的系数（与
%  (3.21) 在最简情形 $\tau_1=B_{12}=0$ 下的 $\tg_{11}=1/R^2$ 也一致）；若把 $E^{ij}$ 误抄成
%  $E_{ij}$，这一核对立刻失败。]
% [Tutor Note: p.13 式 (3.14) 第二式的 $\overline{\phi}$ 不是笔误，可用 (3.5) 与 p.10 的路径积分
%  结论 $J_\mu=\partial_\mu\tX$ 验证：由 $J_0=\partial_1X$、$J_1=\partial_0X$ 得
%  $\partial_0\tX=\partial_1X$、$\partial_1\tX=\partial_0X$，于是
%  $(\partial_0+\partial_1)\tphi=\partial_+\Re\phi+\ii(\partial_0X+\partial_1X)=(\partial_0+\partial_1)\phi$，
%  而 $(\partial_0-\partial_1)\tphi=\partial_-\Re\phi+\ii(\partial_1X-\partial_0X)=(\partial_0-\partial_1)\overline{\phi}$。
%  正是这个“一半对上 $\phi$、另一半对上 $\overline{\phi}$”的不对称，导致 (3.15)–(3.17) 中 $\tphi$
%  由 B-手征变为 A-手征。]
% [Tutor Note: p.13 式 (3.18) 的 bar 结构看似不对称，其实是把 $Q\leftrightarrow\overline{Q}$ 的自同构
%  作用在 $Q_+\to Q'_+$、$Q_-\to\overline{Q}'_-$ 上得到的，故第 3、4 项恰为第 1、2 项的 bar 变换，
%  不是笔误。核对办法：对 (3.18) 的前两式同时作 $Q\leftrightarrow\overline{Q}$（含带撇者），即还原
%  后两式。]
% [Tutor Note: p.14 式 (3.21)–(3.22) 互相印证：把 (3.21) 中 $\tB$ 的系数 $\tau_1/2\pi$ 代入 (3.22)
%  第二式，得 $2\pi\cdot\tau_1/(2\pi)=\tau_1$ 与 $\ii(1/R^2)R^2\tau_2=\ii\tau_2$，相加正是 $\tau$；
%  若 (3.21) 的系数误作 $\tau_2/2\pi$，则无法还原 $\tau$，可据此判断抄写无误。另注意 (3.21) 的度规
%  分量是 $\dd\widetilde{s}^{\,2}$（$s$ 上带波浪号），不是 $\dd s^2$。]
% [Tutor Note: p.14 术语首现：“Lagrangian torus fibration”译作“拉格朗日环面纤维化”（原文用斜体）；
%  “holonomy”译作“和乐”（p.11 首现）；“boundary particle”译作“边界粒子”（p.11）。]
